\section{Bonus: every finite division ring is a field}
\label{sec:wedderburn}

The field axioms require multiplication to commute. Could a finite
arithmetic allow division by every nonzero element while having $ab\ne ba$?
Wedderburn's theorem says no. This is a bonus beyond the finite-fields
lectures. The proof below uses the cyclotomic argument of
\href{https://nonagon.org/ExLibris/sites/default/files/pdf/Uber-Kommutativitat-endlicher-Schiefkorper-Ernst-Witt.pdf}{Ernst Witt};
see also Andr\'e Weil, \emph{Basic Number Theory}, Chapter I, \S1.
We develop the two extra ingredients before using them.

\subsection*{Conjugation partitions a group}

For an element $a$ of a group $G$, its \emph{conjugacy class} is
$\{gag^{-1}:g\in G\}$, and its \emph{centralizer} is the subgroup
\[
C_G(a)=\{g\in G:ga=ag\}.
\]
Two elements $g,h$ give the same conjugate exactly when
$h^{-1}g\in C_G(a)$. Thus the class is in bijection with the left
cosets of $C_G(a)$, so its size is $|G|/|C_G(a)|$.
The \emph{center} $Z(G)$ consists of elements commuting with all of $G$;
these are exactly the one-element conjugacy classes. Since conjugacy
classes partition $G$, choosing one $a_i$ from each noncentral class gives
the \emph{class equation}
\[
|G|=|Z(G)|+\sum_i\frac{|G|}{|C_G(a_i)|}.
\]
This is another use of counting by equal-sized cosets.

\subsection*{The integer factor that will not fit}

Among the complex $n$th roots of $1$, call a root \emph{primitive} if
its multiplicative order is exactly $n$. They are the points
$\exp(2\pi i k/n)$ with $\gcd(k,n)=1$, so there are $\varphi(n)$ of them.
Define the cyclotomic polynomial
\[
\Phi_n(X)=\prod_{\zeta\text{ primitive of order }n}(X-\zeta).
\]
For example $\Phi_1=X-1$, $\Phi_2=X+1$,
$\Phi_3=X^2+X+1$, and $\Phi_4=X^2+1$.
Only the elementary geometry of roots of unity is needed here;
irreducibility of cyclotomic polynomials is not needed.

\begin{lemma}
\label{lem:wedderburn-cyclotomic}
The polynomials $\Phi_n$ are monic with integer coefficients and
\[
X^n-1=\prod_{d\mid n}\Phi_d(X).
\]
If $d$ is a proper positive divisor of $n$, then
\[
\Phi_n(X)\divides\frac{X^n-1}{X^d-1}\quad\text{in }\Z[X].
\]
For integers $q\ge2$ and $n>1$, $|\Phi_n(q)|>q-1$.
\end{lemma}
\begin{proof}
Every $n$th root has a unique order $d\mid n$, giving the product
identity over the complex numbers. To obtain integer coefficients,
induct on $n$, starting with $\Phi_1=X-1$. The product over proper
divisors is monic with integer coefficients by induction. Dividing
$X^n-1$ by it by monic long division produces an integer quotient
and remainder. Over the complex numbers the division is exact with
quotient $\Phi_n$, so the remainder is zero and $\Phi_n\in\Z[X]$.

If $d\mid n$, cancel the factors indexed by divisors of $d$ from the
product for $X^n-1$. When $d<n$, the factor $\Phi_n$ remains, proving
the polynomial divisibility and hence its integer-valued version at $q$.

For the inequality, each primitive $n$th root $\zeta$ is different from
$1$ when $n>1$, and lies on the unit circle. Consequently
\[
|q-\zeta|^2=q^2-2q\operatorname{Re}(\zeta)+1>(q-1)^2.
\]
Taking the product of these distances gives
$|\Phi_n(q)|>(q-1)^{\varphi(n)}\ge q-1$, since $q-1\ge1$
and $\varphi(n)\ge1$.
\end{proof}

\begin{theorem}[Wedderburn's little theorem]
\label{thm:little-wedderburn}
Every finite division ring is commutative, and therefore is a field.
\end{theorem}
\begin{proof}
Let $D$ be a finite division ring and let
\[
Z=\{z\in D:zx=xz\text{ for every }x\in D\}
\]
be its center. This is a field: it is closed under the ring operations,
its multiplication commutes, and the inverse of a central nonzero
element is central, as follows by multiplying $zx=xz$ by $z^{-1}$
on both sides. Write $|Z|=q\ge2$. Addition and multiplication by
central scalars make $D$ a vector space over $Z$; a finite basis exists
by the earlier basis-building argument. If its dimension is $n$, then
$|D|=q^n$. We will show $n=1$.

For $a\in D$, put
\[
C_D(a)=\{x\in D:xa=ax\}.
\]
This set contains $Z$ and is a vector subspace over $Z$, so it has
$q^d$ elements for some $d\ge1$. It is closed under multiplication
and under inverses of nonzero elements: if $xa=ax$, multiplying by
$x^{-1}$ shows $x^{-1}a=ax^{-1}$. Thus its nonzero elements form
a subgroup of $D^\times$ of order $q^d-1$. We do not assume that
this centralizer is commutative. If $a$ is noncentral, the subspace
is proper and $d<n$.

Lagrange gives $q^d-1\mid q^n-1$. This implies $d\mid n$:
write $n=kd+r$, $0\le r<d$. Reduction modulo $q^d-1$ gives
$q^n-1\equiv q^r-1$. If $r>0$, this would be a positive multiple
of $q^d-1$ strictly smaller than $q^d-1$, impossible. Thus $r=0$.

Apply the class equation to $G=D^\times$. Its center is $Z\setminus\{0\}$,
because commuting with every nonzero element is the same as commuting
with all of $D$. Each group centralizer is
$C_D(a)\setminus\{0\}$. Therefore
\[
q^n-1=q-1+\sum_i\frac{q^n-1}{q^{d_i}-1},
\qquad d_i\mid n,\quad d_i<n,
\]
where the sum runs over one representative of each noncentral
conjugacy class.

Suppose $n>1$. By Lemma~\ref{lem:wedderburn-cyclotomic}, the integer
$\Phi_n(q)$ divides $q^n-1$ and every term in this sum. Hence it
divides $q-1$. But $|\Phi_n(q)|>q-1>0$, which is impossible for a
divisor of $q-1$. It follows that $n=1$, so $D=Z$ is a field.
\end{proof}

\begin{bigidea}
The contradiction joins linear algebra, group theory, and integer
divisibility. Dimensions turn centralizers into powers of $q$; conjugation
turns their sizes into a sum; cyclotomic factors divide every large term
but are too large to divide the leftover $q-1$.
\end{bigidea}

\begin{corollary}
A finite ring with $1\ne0$ and no zero divisors is a field, even if
commutativity was not assumed. Thus every finite division ring is,
up to isomorphism, one of the Galois fields $\F_{p^n}$.
\end{corollary}
\begin{proof}
For nonzero $a$, both maps $x\mapsto ax$ and $x\mapsto xa$ are
injective and hence surjective. Choose $b,c$ with $ab=1$ and $ca=1$.
Then $c=c(ab)=(ca)b=b$, so each nonzero element has a two-sided inverse.
The ring is a division ring, to which Wedderburn applies. The finite-field
classification gives the last assertion.
\end{proof}

Finiteness matters. The real quaternions form an infinite division ring
with $ij=k$ and $ji=-k$. Conversely, a finite matrix ring need not be a
division ring: in $M_2(\F_2)$ the nonzero matrices
$\operatorname{diag}(1,0)$ and $\operatorname{diag}(0,1)$ have product zero.
Neither example contradicts the theorem.
